\section{Directed graphs and sparse-base routing}
\label{sec:directed-routes}

\subsection{The directed marker identity}

Assume \HTC. Let a loopless digraph have selected vertices $U$,
selected arcs $F$ with endpoints in $U$, and markers $M_0\subseteq U$.
Replace~\eqref{eq:marker-matrix} by the outgoing Laplacian
\begin{equation}\label{eq:directed-matrix}
 A^+_{vv}=d^+_F(v)+[v\in M_0]+[v\notin U],\qquad
 A^+_{vw}=-[v\mathbin{\to}w\in F]\quad(v\ne w).
\end{equation}
Its determinant counts spanning arborescences directed toward the
new root $\rho$, after adding $v\to\rho$ for every marker $v$.
Here every selected vertex has one outgoing tree arc and every orbit
reaches $\rho$.

We verify this identity directly. Remove the identity block outside
$U$. Expand each diagonal entry as a choice of one outgoing selected
arc or an available root arc. Fix all these choices, a function
$f:U\to U\cup\{\rho\}$. Its nonroot directed cycles are disjoint.
A nontrivial permutation cycle in a determinant term must be one of
these cycles: its off-diagonal entries require the corresponding
choices of $f$. Conversely any subset of these cycles may be used,
with every other vertex a diagonal fixed point. A length-$d$ cycle
has permutation sign $(-1)^{d-1}$ and $d$ negative entries, so
contributes $-1$. The total for this $f$ is
\[
 \sum_{Q\subseteq\operatorname{cycles}(f)}(-1)^{|Q|}
 =(1-1)^{|\operatorname{cycles}(f)|}.
\]
It is one exactly when every orbit reaches the root. Thus
$\det A^+\ge0$, and it is positive exactly when every selected vertex
can reach a marker along selected arcs. For the converse existence
claim, choose for each vertex an arc decreasing its shortest distance
to a marker or to the new root. Empty $U$ has determinant one.
The exclusion of loops keeps diagonal choices separate from
permutation cycles.

The potential matrix positions number at most $n+m$, and outgoing
degree counters and entry scopes are $O(n+m)$. Combining this
identity with the sparse sign gadget obeys the same wire, degree,
entry-cost, and ring bounds as before.

\subsection{Directed applications}

\begin{proof}[Proof of Theorem~\ref{thm:directed}]
For directed HC, require indegree and outdegree one at every vertex
and fix one marker $r$. Candidates are disjoint directed cycles.
A spanning cycle has exactly one rooted arborescence in the augmented
graph: replace $r$'s outgoing cycle arc by its root arc, and keep the
other arcs. Any additional cycle misses the marker and has determinant
zero. The final sum is therefore the directed HC count itself; it
is not divided by $n$.

For a directed $s$--$t$ path, $s\ne t$, require selected indegree/outdegree
$(0,1)$ at $s$, $(1,0)$ at $t$, $(1,1)$ at other selected vertices,
and $(0,0)$ outside $U$. Use $t$ as the sole marker. The path contributes
one, and any extra directed cycle makes the determinant zero. This
proves exact counting with a prescribed length and longest-path
optimization. Unspecified endpoints are enumerated as \emph{ordered}
pairs. A single-vertex path is treated separately. A two-vertex directed
cycle with opposite arcs is a valid HC under these conventions.
Empty graphs and the one-vertex loopless case are checked directly.

For directed Steiner tree toward a root $r$, require $r\in U$,
outdegree zero at $r$, outdegree one at each other selected vertex,
and inclusion of all terminals. With marker $r$ the determinant is
one precisely on arborescences toward $r$. A feasible rooted-reachability
subgraph contains a suitable arborescence, obtained by selecting arcs
that decrease distance to the root; extra branches not needed for the
terminals can be deleted. With nonnegative costs this restriction
preserves the optimum. Reverse the input for the usual version with
paths from the root to the terminals. The same network exactly counts
arborescent subgraphs of the specified direction. Even when terminals
consist only of the root, the counting problem includes all such
subgraphs; only the optimum can immediately be the single root.

For a strongly connected selected subgraph, fix $r\in U$ and multiply
two determinant gadgets on the \emph{same} selected-arc bits: one for
$F$ and one for its reverse. The second matrix's diagonal uses the
original indegrees; it is recomputed rather than obtained by merely
transposing the original matrix. The product is positive exactly when
every selected vertex can reach $r$ and can be reached from $r$.
This is strong connectivity. The product stays inside the common
sum over candidates. It is not the product of two separately
contracted sums. For a spanning subgraph use $U=V$. For a terminal
version use $T\subseteq U$ and a fixed $r\in T$ when $T$ is nonempty.
For an empty terminal set use the requested empty-subgraph convention,
or enumerate roots to cover nonempty solutions.

Each arborescence is a subset of at most $m+n$ potential augmented
arcs. At most two determinant factors and at most $2^{n+m}$
candidates give the uniform bound
\[
 0\le Z\le2^{3(n+m)}<2^{4(n+m+2)}.
\]
Use $P_0=4(n+m+2)+1$, $\ell=1$, and a fixed multiple of $n+m+2$
for the \HTC\ scale. Costs use the exact statistics gadget; objective
and witness bits are recovered by the same positive-sum self-reduction
as in the undirected case. All query counts are polynomial. This
proves the theorem.
\end{proof}

\subsection{Closed walks, required edges, and metric TSP}

\begin{proof}[Proof of Theorem~\ref{thm:routing}]
A nonempty closed walk determines a connected even-degree edge
multiset. If an edge has multiplicity at least three, delete two of
its copies. Its support remains unchanged, degrees stay even, and
required-edge and terminal constraints remain satisfied. The cost
does not increase. Repeating gives an optimum with multiplicities
\begin{equation}\label{eq:walk-multiplicity}
 \mu_e\in\{0,1,2\}.
\end{equation}
Conversely a connected even-degree multigraph has an Euler tour:
follow unused edges to a closed trail, and splice in further closed
trails wherever unused edges remain. Connectivity and parity ensure
this process exhausts every edge. A selected single vertex without
edges represents its zero-length closed walk.

Encode $\mu_e=a_e+2b_e$ with Boolean bits and forbid $a_e=b_e=1$.
Define its support $f_e=a_e\vee b_e$ uniquely. Required edges have
$f_e=1$. Select $U$ containing the terminals and every support
endpoint; the single-marker determinant on the support graph
$(U,\{e:f_e=1\})$ enforces connectivity. Extra isolated selected
vertices are thereby eliminated except for the valid singleton case.
For nonempty requirements fix
$r\in T\cup V(R)$ as the marker and require $r\in U$. If a depot is
specified, include it in $T$. Empty requirements allow an empty walk
of cost zero and are handled directly.

A parity chain at vertex $v$ requires
\[
 \sum_{e\ni v}a_e\equiv0\pmod2,
\]
which is exactly the parity of its multidegree. The chains and the
multiplicity/support relations together have $O(n+m)$ wires. The
cost is the exact statistic
\[
 \sum_e w_ea_e+2w_eb_e.
\]
Its coefficient width is at most one more than the input width;
this constant change is covered by the fixed-scale long-coefficient
gadget. The determinant counts trees of the \emph{simple support}
graph, as a positive feasibility weight.

There are at most $2^n3^m$ candidates $(U,\mu)$, and at most $2^m$
support spanning trees for a single marker. Thus
\[
 0\le Z\le2^n6^m<2^{4(n+m+2)}.
\]
The same precision and \HTC\ scale apply. Binary search on cost,
followed by fixing $a_e,b_e,u_v$ while preserving positive sums,
recovers an optimal even connected multiset. Its Euler tour has at
most $2m$ edge traversals. With a depot, rotate the tour to start
and end there.

Finally suppose the base graph is connected, $T=V$, and $R$ is empty.
Expand any Hamiltonian tour of its shortest-path metric into base-graph
shortest paths. This gives a closed spanning walk of the same cost.
Conversely, visit vertices in their order of earliest appearance along a
spanning walk and shortcut repeats in the metric. The triangle
inequality gives a Hamiltonian metric tour of no greater cost.
Therefore the optima are equal, and shortcutting an optimal returned
walk gives an optimal metric tour. Zero costs can give a pseudometric;
the same inequalities still hold. Distances are needed only for this
polynomial-time output step; a dense metric closure is not part of
the tensor network. The proof of~\eqref{eq:graph-time} follows.
\end{proof}
